%%%%%%%%%%%%%%%%%%%%%%%%%%%%%%%%%%%%%%%%%%%%%%%%
% Mülleimer (Schild über den Mülleimern)
% CC-BY-SA 3.0, FAU FabLab
% Früher Hinweisschilder_divers/muelleimer.pdf (Illustrator), jetzt LaTeX.
%%%%%%%%%%%%%%%%%%%%%%%%%%%%%%%%%%%%%%%%%%%%%%%%
\documentclass[a4paper,landscape]{scrartcl}
\usepackage{fablab-aushang}
\Fabcube

\begin{document}
\Schild
\SchildKopf{Mülleimer \en{Trash}}

\begin{Spalte}[0.68\linewidth]
\begin{itemize}
  \item Volle Müllsäcke zuknoten und vor die Tür stellen.\\
        \en{Place full bags in front of the door.}
  \item Neue Müllsäcke liegen unten.\\
        \en{Trash bags are underneath.}
\end{itemize}
\end{Spalte}

\vfill
% Die Pfeile zeigen auf die beiden Mülleimer darunter
\begin{tikzpicture}[x=\linewidth, y=\SchildText, baseline=(current bounding box.south)]
  \node[fill=fcBlau, rounded corners=0.3\SchildText, minimum width=0.46\linewidth,
    minimum height=2.2\SchildText, font=\GrossSchrift] (pa) at (0.24,0) {Papier \en{Paper}};
  \node[fill=fcRot, rounded corners=0.3\SchildText, minimum width=0.46\linewidth,
    minimum height=2.2\SchildText, font=\GrossSchrift] (re) at (0.76,0) {Restmüll \en{Landfill}};
  \draw[-{Triangle[length=40pt, width=40pt]}, line width=11pt]
    ([yshift=-0.4\SchildText]pa.south) -- ++(-0.07,-3.4);
  \draw[-{Triangle[length=40pt, width=40pt]}, line width=11pt]
    ([yshift=-0.4\SchildText]re.south) -- ++(0.07,-3.4);
\end{tikzpicture}

\end{document}
